% Part: counterfactuals
% Chapter: introduction
% Section: counterfactuals

\documentclass[../../../include/open-logic-section]{subfiles}

\begin{document}

\olfileid{cnt}{int}{cnt}

\olsection{প্রতিবাস্তব শর্তবচন}

ইংরেজির ``if \dots then \dots'' গঠনের একটি প্রচলিত
ও গুরুত্বপূর্ণ রূপে \emph{to be} ক্রিয়ার অতীত সম্ভাবক
রূপ ব্যবহৃত হয়: ``যদি \dots ঘটত, তবে \dots ঘটত''।
এ ধরনের শর্তবচনের পূর্বপক্ষ সাধারণত মিথ্যা—অর্থাৎ বাস্তব
ঘটনার বিপরীত—বলে এগুলিকে \emph{প্রতিবাস্তব শর্তবচন}
বলা হয়। \emph{to be} ক্রিয়ার সম্ভাবক রূপ ব্যবহৃত হয়
বলে এগুলির আরেক নাম \emph{সম্ভাবক-ভাবের শর্তবচন}।
এগুলি \emph{নির্দেশক} শর্তবচন থেকে আলাদা; নির্দেশক
শর্তবচনের রূপ ``যদি \dots ঘটে, তবে \dots ঘটে''।
দুই ধরনের শর্তবচনের সত্যতার শর্ত ভিন্ন। অ্যাডামসের
বিখ্যাত উদাহরণটি দেখি:
\begin{quote}
  অসওয়াল্ড কেনেডিকে হত্যা না করে থাকলে অন্য কেউ করেছে।

  অসওয়াল্ড কেনেডিকে হত্যা না করলে অন্য কেউ করত।
\end{quote}
প্রথমটি নির্দেশক, দ্বিতীয়টি প্রতিবাস্তব শর্তবচন। প্রথমটি
স্পষ্টত সত্য: আমরা জানি মার্কিন প্রেসিডেন্ট জন এফ.
কেনেডিকে \emph{কেউ} হত্যা করেছিল। ওয়ারেন কমিশনের
প্রতিবেদনের বিপরীতে সেই ব্যক্তি লি হার্ভি অসওয়াল্ড
না হলে অন্য কেউ তাঁকে হত্যা করেছিল। দ্বিতীয় বাক্যটির
দাবি আলাদা। এতে বলা হচ্ছে, অসওয়াল্ড কেনেডিকে হত্যা
না করলে—অর্থাৎ ডালাসে গুলি চালনার ঘটনাটি এড়ানো গেলে
বা তা ব্যর্থ হলে—পরবর্তী ইতিহাস এমনভাবে এগোত যে
আরেকটি হত্যাচেষ্টা সফল হতো।
% BN-SRC-764: ষড়যন্ত্র এই প্রতিবাস্তব দাবির একমাত্র আবশ্যক শর্ত নয়; অন্য সফল হত্যাচেষ্টার অনুমান পৃথক।
উৎসের এই ব্যাখ্যায় কেনেডির মৃত্যু নিশ্চিত করার ষড়যন্ত্র
একটি সম্ভাব্য প্রেক্ষাপট, যেমন বহু ষড়যন্ত্রতত্ত্বে বিশ্বাসী ব্যক্তি
মনে করেন। তবে স্বতন্ত্র অন্য হত্যাকারীর সফল চেষ্টা থেকেও
বাক্যটি সত্য হতে পারে; ষড়যন্ত্রই তার একমাত্র আবশ্যক শর্ত নয়।

\emph{নির্দেশক} শর্তবচনের সত্যতার শর্ত বস্তুগত
শর্তবচন যথাযথভাবে প্রকাশ করে কি না, তা নিয়ে এখনো
বিতর্ক আছে। বিশেষত, বস্তুগত শর্তবচনের আপাত-বৈপরীত্য
এমনভাবে ``ব্যাখ্যা'' করা যায় কি না, যাতে নির্দেশক
শর্তবচনের সত্যতার শর্ত বস্তুগত শর্তবচন দিয়েই ধরা যায়,
তা বিতর্কিত। কিন্তু প্রতিবাস্তব শর্তবচনকে বস্তুগত
শর্তবচন দিয়ে যথাযথভাবে প্রতীকায়িত করা যায় না। কারণ
প্রতিবাস্তব শর্তবচনের পূর্বপক্ষ সাধারণত মিথ্যা হলেও,
মিথ্যা পূর্বপক্ষযুক্ত সব প্রতিবাস্তব শর্তবচন সত্য নয়।
যেমন, ওয়ারেন কমিশনের প্রতিবেদন সত্য এবং কেনেডিকে
হত্যার কোনো ষড়যন্ত্র ছিল না বলে বিশ্বাস করার পাশাপাশি,
অসওয়াল্ডের চেষ্টা ব্যর্থ হলে স্বতন্ত্র অন্য সফল হত্যাচেষ্টাও
ঘটত না বলে ধরে নিলে অ্যাডামসের প্রতিবাস্তব শর্তবচনটি
মিথ্যা বলে ধরা যায়।

কারণসংক্রান্ত যুক্তিবিচারে প্রতিবাস্তব শর্তবচনের গুরুত্বপূর্ণ
ভূমিকা আছে: কার্যকারণ সম্পর্ক প্রকাশ তার একটি প্রধান
ব্যবহার। যেমন, দেশলাই-কাঠি ঘষলে সেটি জ্বলে ওঠে;
``এই কাঠিটি ঘষা হলে জ্বলে উঠত'' বলে সম্পর্কটি প্রকাশ
করা যায়। বস্তুগত, এমনকি সাধারণভাবে নির্দেশক,
শর্তবচন দিয়ে তা প্রকাশ করা যায় না। ``কাঠিটি ঘষা হয়
$\lif$ কাঠিটি জ্বলে ওঠে''—কাঠিটি কখনো ঘষা না হলে
বক্তব্যটি সত্য, ঘষা হলে কী ঘটত তা নির্বিশেষে। আরও
সমস্যা হলো, ``কাঠিটি ঘষা হয় $\lif$ কাঠিটি ফুলের
তোড়ায় পরিণত হয়''—কাঠিটি ঘষা না হলে এটিও সত্য;
অথচ ঘষা হলে কাঠিটি নিশ্চয়ই ফুলের তোড়ায় পরিণত হতো না।

প্রতিবাস্তব শর্তবচনের সঠিক যুক্তিবিদ্যা ঠিক কী, তা নিয়ে
এখনো বিতর্ক আছে। স্টালনেকার ও লুইস একটি প্রভাবশালী
বিশ্লেষণ দিয়েছেন। তাঁদের মতে ``যদি~$S$ ঘটত, তবে~$T$
ঘটত'' প্রতিবাস্তব শর্তবচনটি সত্য তখন এবং কেবল তখনই, যখন $S$ সত্য
এমন সম্ভাব্য পরিস্থিতিগুলির মধ্যে বাস্তব জগতের সঙ্গে
সবচেয়ে নিকটবর্তী পরিস্থিতিতে (``সম্ভাব্য জগতে'') $T$
সত্য।
% BN-SRC-765: এক নির্বাচিত নিকটতম জগৎ ও সব সমনিকট জগৎ পৃথক; লুইসের সাধারণ অর্থতত্ত্বে নিকটতম জগৎ থাকা বাধ্যতামূলক নয়।
এই কথাটি নিকটতম জগতের একটি প্রাথমিক সারসংক্ষেপ।
স্টালনেকারের বিশ্লেষণে এক নির্বাচিত নিকটতম জগৎ ব্যবহৃত হয়;
লুইসের ক্ষেত্রে নিকটতম জগৎ থাকলে সেই সব জগতে অনুপক্ষ
সত্য হতে হয়। তাঁর সাধারণ বিশ্লেষণ নিকটতম জগৎ না-থাকার
ক্ষেত্রও সামলায়; পরের অধ্যায়ে গোলকের শর্তে সেই সংজ্ঞা দেওয়া
হয়। এই বিশ্লেষণকে ``সত্তাতাত্ত্বিক'' বলা হয়, কারণ
এতে সম্ভাব্য জগতের অস্তিত্বতত্ত্বের উল্লেখ আছে। অন্য
বিশ্লেষণগুলি শর্তাধীন সম্ভাবনা কিংবা বিশ্বাস-সংশোধনের
তত্ত্ব ব্যবহার করে। প্রতিবাস্তব শর্তবচনের জন্য নানা
যুক্তিবিদ্যা প্রস্তাবিত হয়েছে। এমনকি একটিমাত্র
লুইস--স্টালনেকার যুক্তিবিদ্যাও নেই: নিকটতম সম্ভাব্য
জগতের ধারণায় দুজনের ব্যাখ্যা কাছাকাছি হলেও, তাঁদের
ভিন্ন অনুমান থেকে ভিন্ন বৈধ অনুসিদ্ধান্ত পাওয়া যায়।

\end{document}
